%----------------------------------------------------------------------------
\chapter{Archiving selected research artifacts}\label{guide:zenodo}
%----------------------------------------------------------------------------
A reviewer assessing \textbf{correctness} (criterion D/2) needs the artifacts behind your results. Archive the selected code, proofs and benchmarks on \href{https://zenodo.org/}{Zenodo} and cite the verified record. The project monorepo stays private. Its \path{code/}, \path{proof/} and \path{benchmark/} directories are ordinary components, not separate repositories or submodules.

The lead student owns the Zenodo records and publishes after advisor review. This chapter follows the repository's manual preparation and publication workflows. Paths below are relative to the monorepo root; \path{scripts/release/README.md} documents the commands and recovery procedures. Zenodo's \href{https://developers.zenodo.org/}{API documentation} describes its metadata and upload operations.

\section{Select the files and metadata}
\begin{enumerate}
  \item Complete the reproduction README using \cref{guide:code:readme}. Enable the actual checks for each selected component in \path{project.json}; an unconfigured scaffold is not evidence that research results reproduce.
  \item Create a tracked selection manifest, normally \path{release/manifest.json}, using the examples in \path{docs/selective-release/}. Replace their example paths with the project's real files. Select the code, proof obligations, benchmark inputs/results and documentation needed to reproduce the claims. Include the Dockerfile and build inputs when they are part of that procedure. A separately distributed container image needs its own publication plan and recorded image/platform digest.
  \item Choose the archive's licenses and access conditions with the advisors. Preserve third-party notices and redistribution terms. The default is one combined code/proof/benchmark record; use independently citable records only when the project needs them.
  \item Set each record's \texttt{metadata\_file} in the manifest, for example \path{release/metadata/study.json}. Fill its title, creators, version, date, access and license (\cref{guide:lst:zenodo-json}). List the actual supervisors as contributors. Keep any selected \texttt{CITATION.cff} consistent with these fields. The manifest's access setting must agree with the metadata.
  \item Add real related identifiers when available (\cref{guide:lst:zenodo-related}). Credit reused work and retain its license; leave an unknown identifier out instead of inventing a DOI. Replace every illustrative placeholder before preparing a candidate.
\end{enumerate}

\begin{listing}[!ht]
\begin{minted}{json}
{
  "title": "TODO Artifact for: <thesis title>",
  "description": "Code, proofs and benchmarks for <thesis title>.",
  "upload_type": "software",
  "version": "1.0.0",
  "publication_date": "TODO YYYY-MM-DD",
  "access_right": "open",
  "license": "MIT",
  "creators": [
    { "name": "TODO Surname, Given name", "affiliation": "Budapest University of Technology and Economics", "orcid": "TODO your ORCID iD, or remove this field" }
  ],
  "contributors": [
    { "name": "Farkas, Martin", "type": "Supervisor", "affiliation": "Budapest University of Technology and Economics" },
    { "name": "Kocsis, Imre", "type": "Supervisor", "affiliation": "Budapest University of Technology and Economics" }
  ],
  "keywords": ["TODO"]
}
\end{minted}
\caption{Illustrative record metadata in the file selected by the manifest. Confirm the actual creators, supervisors, license and access; this example is not ready to publish.}\label{guide:lst:zenodo-json}
\end{listing}

\begin{listing}[!ht]
\begin{minted}{json}
"related_identifiers": [
  { "identifier": "https://github.com/<upstream>/<repo>", "relation": "isDerivedFrom", "resource_type": "software" },
  { "identifier": "10.xxxx/<paper DOI>", "relation": "isSupplementTo", "resource_type": "publication-article" }
]
\end{minted}
\caption{Related works in record metadata. Replace these example identifiers with verified public identifiers, or omit them.}\label{guide:lst:zenodo-related}
\end{listing}

\section{Review what the archive discloses}\label{guide:zenodo:public}
The packager exports the selected bytes and their provenance. It does not export the whole Git history. Making the GitHub repository public is a separate decision with a separate history review; publication of selected artifacts does not require it.
\begin{itemize}
  \item Inspect every selected file, its metadata and the generated provenance. Keep \path{scratch/}, private integration links, internal review records, account state and agent configuration out of the archive unless explicitly selected and reviewed. A private source URL or an author field can disclose information even when the research files have restricted access.
  \item Publish only the paper material that the manifest deliberately selects. Selected paper files require strict submission evidence; ordinary advisory writing CI is not sufficient for that release.
  \item Keep secrets and test keys in ignored local state. Rotate a credential that was committed; deleting a current file does not remove it from Git history. Use relative reproduction paths rather than machine-specific paths.
  \item Check that the README explains the selected files, commands, limits and results mapping. Review third-party attribution and redistribution rights for code, data, fonts and artwork. Hashes identify bytes; they do not establish permission to publish them.
\end{itemize}
Use \path{code/scripts/publishing-checklist.md} to record this review, optionally in a private \path{code/PUBLISHING.md}. The checklist records advisor review; creating it does not enable a workflow. The code hygiene checker is a supplementary local check, not the release gate. Student-writable workflows do not enforce advisor approval, so the student and advisor must complete that review before publication.

\section{Prepare and publish a reviewed candidate}
Before the first production publication, complete the sandbox pilot with the advisors and check the project's immutable-release setting. The repository administrator provisions the lead student's own \texttt{ZENODO\_SANDBOX\_TOKEN} and \texttt{ZENODO\_TOKEN} as project secrets in GitHub's repository Actions-secret settings. Tokens belong in secret storage, never in metadata, manifests or command-line arguments.
\begin{enumerate}
  \item Commit the manifest, metadata and research source, and obtain successful checks for that exact commit and configuration. Retain the checks run ID, attempt, artifact ID and archive digest. The selected artifacts must describe the measurements and proof obligations actually reported in the paper.
  \item From the default branch, manually run \emph{Prepare research release} (\path{.github/workflows/release-prepare.yml}). Supply that source SHA, checks run/attempt/artifact identity and digest, and the manifest path. It packages verified existing bytes without rebuilding research code.
  \item Download and inspect the candidate ZIPs, file inventory, metadata, runtime evidence and \texttt{receipt.json}. Record the candidate's preparation run, attempt, artifact ID, archive digest and exact receipt digest from the workflow summary. The advisor reviews these particular bytes; a successful build alone is not publication approval.
  \item Manually run \emph{Publish reviewed research} (\path{.github/workflows/publication.yml}) from the default branch. Supply the candidate identity, reviewed receipt digest and a stable operation ID. Choose \texttt{sandbox} for the pilot, or deliberately choose \texttt{production} for a reviewed real release. Start with phase \texttt{stage} to prepare the Zenodo draft. Before any Zenodo write, the workflow retains the verified candidate in an immutable project support release so Actions artifact expiry cannot silently replace it.
  \item Inspect the draft and, after advisor review, invoke the same operation and candidate with phase \texttt{publish}. Staging or reserving a DOI does not make the archive available. Only a receipt verifying the production record as \texttt{published}, with matching files and metadata, supports an availability statement or DOI badge.
  \item If a response is interrupted or uncertain, retain the same operation and candidate identity and use \texttt{reconcile}. Follow the recovery instructions in \path{scripts/release/README.md} if it remains blocked. Creating another operation to bypass an uncertain write risks duplicate records.
\end{enumerate}
The journal on the project's \texttt{publication-state} branch preserves retry state. It is excluded from exported research packages and template contributions. GitHub tags, support releases and ordinary pushes do not automatically publish to Zenodo; use the manual workflow above.

\section{Cite it}\label{guide:zenodo:cite}
Zenodo's \href{https://support.zenodo.org/help/en-gb/1-upload-deposit/97-what-is-doi-versioning}{versioning documentation} distinguishes a \emph{version DOI} for one archived version from a \emph{concept DOI} for all its versions.
\begin{itemize}
  \item Cite the concept DOI for the artifact as a whole. Beside the results, identify the exact version DOI and \textbf{results tag}: the Git tag and source commit used for the reported numbers. A later packaging version must not change which version those measurements describe.
  \item Generate the README badge, version citation and artifact-availability patch from the verified production publication receipt using \path{scripts/release/writeback.py}. Review that patch before applying it. The bibliography target is \path{paper/bib/references.bib}; preserve Zotero's managed \path{paper/bib/zotero.bib}.
  \item Use the version and concept DOIs returned by the verified record. Draft/reserved identifiers and sandbox records are not evidence of production availability, and DOI numbers must never be guessed.
\end{itemize}

\section{Release a new version}\label{guide:zenodo:versions}
For a new version, update the selected metadata, commit the changed source, run the required checks and prepare a new reviewed candidate. Give that candidate a new operation ID and repeat the stage/review/publish sequence. The project journal links it to the prior record; preserve that journal and the earlier receipts. Keep old tags and archived bytes fixed.

The lifecycle (\cref{guide:tab:lifecycle}) allows code polish after the paper deadline. Within about two days of upload, aim to complete the README, container reproduction instructions, results map and reproduction limits, then publish that polish as a new version. Keep the paper's results tag and measured version unchanged. Later polish is allowed, but reviewers may already have looked.
